\section{Теоретическая часть}

Антенна --- устройство, предназначенное для излучения или приема волн (в нашем случае --- электромагнитных). Одна из важнейших функций антенны состоит в формировании излучения с определёнными направленными свойствами. Основными характеристиками направленности антенны являются диаграмма направленности (ДН) по амплитуде или по мощности, коэффициент направленного действия (КНД) и коэффициент усиления (КУ). Напомним, как вводятся эти характеристики.

Диаграмма направленности по амплитуде есть угловое распределение амплитуды поля излучения, т.е. зависимость этой амплитуды от полярного $\theta$ и азимутального $\varphi$ углов при фиксированном расстоянии $r$ от антенны. Диаграмма направленности по мощности есть угловое распределение мощности излучения в единицу телесного угла $P(\theta, \varphi) = r^{2} S_{r}(\theta, \varphi)$, где $S_{r}$ --- радиальная компонента вектора Пойнтинга на достаточно большом расстоянии $r$ от антенны. Представляется удобным использование (наряду с абсолютной) нормированной диаграммы направленности $F(\theta, \varphi) = P(\theta, \varphi) / P(\theta_{m}, \varphi_{m})$, где $P(\theta_{m}, \varphi_{m})$ --- мощность, излучаемая в единичный телесный угол в направлении главного максимума $(\theta_{m}, \varphi_{m})$ диаграммы направленности. Диаграммы направленности изображают графически либо в виде «объёмной», рельефной карты, где по каждому угловому направлению $(\theta, \varphi)$ откладывается величина, пропорциональная амплитуде поля излучения или излучаемой мощности (см.~\cref{fig:1a}), либо с помощью плоской проекции отдельных, чаще всего двух ортогональных сечений, проходящих через направление главного максимума и векторы электрического $\mathbf{E}$ и магнитного $\mathbf{H}$ полей (см.~\cref{fig:1b}). Поскольку основная часть мощности, излучаемой направленной антенной, сосредоточена, как правило, в главном лепестке, то весьма показательной представляется его угловая ширина, определяемая обычно по уровню половины мощности $(\Delta\theta_{0.5})$, а иногда и по нулевому (или минимальному) значению $(\Delta\theta_{0})$, как показано на \cref{fig:1b}. Диаграмма направленности антенны, характерный размер $l$ излучающей апертуры которой порядка или больше длины излучаемой волны $\lambda$, окончательно формируется в зоне Фраунгофера, определяемой соотношением
\begin{align} \label{eq:th:1}
	r \gg \frac{l^{2}}{\lambda}.
\end{align}

\begin{figure}[t!]
	\centering
	\begin{minipage}{0.9\textwidth}
		\centering
		\hspace{0.075\linewidth}
		\includegraphics[width=0.4\linewidth]{1a}
		\subcaption{ }
		\label{fig:1a}
	\end{minipage}
	\begin{minipage}{0.9\textwidth}
		\centering
		\includegraphics[width=0.6\linewidth]{1b}
		\subcaption{ }
		\label{fig:1b}
	\end{minipage}	
	\caption{Диаграмма направленности}
	\label{fig:1}
\end{figure}

\textit{Коэффициент направленного действия} $D$ характеризует выигрыш по мощности в направлении максимального излучения вследствие направленности антенны. Он равен отношению мощности, излучаемой в единицу телесного угла в направлении максимума диаграммы направленности $P(\theta_m, \varphi_m)$, к средней мощности $P_{\text{ср}} = P_{\text{изл}}/(4\pi)$, излучаемой антенной по всем направлениям, т.е. $D = 4\pi P(\theta_m, \varphi_m)/P_{\text{изл}}$, где $P_{\text{изл}}$ --- полная излучаемая мощность:
\[
P_{\text{изл}} = \int\limits_{0}^{2\pi} \int\limits_{0}^{\pi} P(\theta, \varphi) \sin \theta \, d\varphi d\theta.
\]
Таким образом, имеем
\begin{align} \label{eq:th:2}
	D = \frac{4\pi P(\theta_m, \varphi_m)}{\int\limits_{0}^{2\pi} \int\limits_{0}^{\pi} P(\theta, \varphi) \sin \theta \, d\varphi d\theta}.
\end{align}

\textit{Коэффициент усиления} $G$ определяется как произведение КНД на коэффициент полезного действия (КПД) антенны $\eta$ (или, точнее, всего антенного тракта):
\begin{align} \label{eq:th:3}
	G = D\eta.
\end{align}
Этот последний коэффициент в свою очередь есть отношение полной мощности $P_{\text{изл}}$, излучаемой антенной, к полной мощности $P_{\text{подв}}$, подводимой к антенне, т.е.
\begin{align} \label{eq:th:4}
	\eta = \frac{P_{\text{изл}}}{P_{\text{подв}}} = \frac{\int\limits_{0}^{2\pi} \int\limits_{0}^{\pi} P(\theta, \varphi) \sin \theta \, d\varphi d\theta}{P_{\text{подв}}}.
\end{align}

В силу принципа взаимности ДН и КНД антенны при ее работе в режиме передачи и в режиме приема совпадают.

Для адекватного описания \textit{приемной антенны} вводятся некоторые дополнительные характеристики. Одна из основных таких характеристик --- эффективная площадь приема антенны $A$.

\textit{Эффективная площадь приема} $A$ определяется как отношение полной принимаемой антенной мощности $P_{\text{пр}}$ к плотности потока падающего излучения $S_{\text{п}}$ в месте расположения антенны:
\begin{align} \label{eq:th:5}
	A = \frac{P_{\text{пр}}}{S_{\text{п}}}.
\end{align}

Как показано в \cref{eq:th:1,eq:th:2}, величины $A$ и $D$ связаны соотношением
\begin{align} \label{eq:th:6}
	A = \frac{\lambda^2}{4\pi} D.
\end{align}

Цель настоящей работы заключается в экспериментальном определении КНД пирамидальной рупорной антенны с помощью так называемого зеркального метода (метода Пареслла) и в сравнении измеренного значения с рассчитанным теоретически. Зеркальный метод опирается на использование идеально (зеркально) отражающей плоской поверхности, расположенной в зоне Фраунгофера и ориентированной параллельно излучающей апертуре (см.~\cref{fig:2}).
\begin{figure}[H]
	\centering
	\includegraphics[width=0.7\linewidth]{2}
	\caption{Метод изображений}
	\label{fig:2}
\end{figure}

Согласно методу изображений отыскание отраженного поля, поступающего в антенну, сводится к нахождению поля, принимаемого от аналогичной зеркальной (относительно отражающей плоскости) излучающей антенны (\cref{fig:2}). В результате последовательного пересчета имеем: мощность, излучаемая гипотетической зеркальной антенной в единицу телесного угла в направлении на реальную антенну, равна $P_{\text{п}} = DP_{\text{изл}}/(4\pi)$, откуда плотность потока энергии в месте приема $S_{\text{п}} = P_{\text{п}}/4X^2 = DP_{\text{изл}}/(16\pi X^2)$, где $X$ --- расстояние между антенной и отражающей плоскостью; наконец, мощность, принимаемая антенной, равна $P_{\text{пр}} = AS_{\text{п}} = ADP_{\text{изл}}/(16\pi X^2)$. С учетом \cref{eq:th:6} окончательно получаем
\begin{align} \label{eq:th:7}
	\frac{P_{\text{пр}}}{P_{\text{изл}}} = \frac{D^2\lambda^2}{64\pi^2 X^2}.
\end{align}

\begin{figure}[t]
	\centering
	\includegraphics[width=0.7\linewidth]{3}
	\caption{Блок-схема установки: 1 --- генератор, 2 --- измерительная линия, 3 --- индикатор, 4 --- согласующее устройство, 5 --- рупорная антенна, 6 --- поглощающий щит, 7 --- отражающий щит}
	\label{fig:3}
\end{figure}
Отсюда интересующая нас величина $D$ представляется в виде
\begin{align} \label{eq:th:8}
	D = \frac{8\pi X}{\lambda} \sqrt{\frac{P_{\text{пр}}}{P_{\text{изл}}}}.
\end{align}

Таким образом, экспериментальное определение КНД требует нахождения отношения принимаемой зеркально отраженной мощности к мощности, излучаемой пирамидальной рупорной антенной.

Измерительная установка включает генератор СВЧ диапазона (длина излучаемой волны $\lambda \approx 3$ см) с отдельным блоком питания, волноводный тракт с измерительной линией и индикатором к ней, пирамидальный рупор, отражающий щит, щит с поглощающим покрытием. Блок-схема установки представлена на \cref{fig:3}. Отражающий щит должен располагаться в зоне Фраунгофера $X \gg l_{1,2}^2/\lambda$ ($l_{1,2}$ --- линейные размеры раскрыва рупора) и иметь линейные размеры $L_{1,2}$, позволяющие перекрывать основную лепестков диаграммы направленности: $L_{1,2} > X\Delta\theta_{1,2} \approx X\lambda/l_{1,2}$ ($\Delta\theta_{1,2}$ --- ширина основного лепестка в горизонтальной или вертикальной плоскости). Убедитесь, что эти условия выполнены! Установка позволяет контролируемо менять расстояние $X + \Delta X$ между антенной и отражающим щитом в пределах $\Delta X \sim 100$ см. Следует особо подчеркнуть, что в измерительной линии используется квадратичный детектор, поэтому показания индикатора пропорциональны квадрату напряженности электрического поля в месте расположения зонда измерительной линии.

В согласованном (с хорошей точностью) режиме, когда отражение от конца волновода отсутствует, коэффициент отражения $\Gamma$ волны в волноводном тракте совпадает с членом $\sqrt{P_{\text{пр}}/P_{\text{изл}}}$, содержащимся в (8), и очевидным образом представляется через коэффициент бегущей волны (КБВ) в волноводе $\kappa = E_{\min}/E_{\max}$ ($\Gamma = (1 - \kappa)/(1 + \kappa)$), определяемый с помощью измерительной линии. Здесь $E_{\min}$ и $E_{\max}$ --- соответственно минимальное и максимальное значения амплитуды поля в волноводном тракте. В результате получаем следующую формулу для КНД:
\begin{align} \label{eq:th:9}
	D = \frac{8\pi X}{\lambda} \frac{1 - \kappa}{1 + \kappa}.
\end{align}

Если детектор в измерительной линии квадратичный, то индикатор дает значения, пропорциональные $|E|^2$, так что вместо $\kappa$ измеряется величина $K = \kappa^2$; при этом вместо (9) нужно использовать выражение
\begin{align} \label{eq:th:10}
	D = \frac{8\pi X}{\lambda} \frac{1 - \sqrt{K}}{1 + \sqrt{K}}.
\end{align}

Итак, при наличии эффективного и надежного согласующего устройства отыскание КНД зеркальным методом сводится к процедуре согласования и последующего измерения КБВ в подводящем волноводном тракте. Согласование достигается за счет включения в волноводный тракт соответствующего устройства (показано пунктиром на \cref{fig:3}) при использовании дополнительного щита с поглощающим покрытием, перехватывающего поле излучения (штрих-пунктирная линия на том же рисунке).

Вы будете работать в несогласованном режиме, когда согласующее устройство отсутствует и специальной процедуры согласования не проводится. С учетом отражения от конца подводящего тракта поле на оси волновода, нормированное на амплитуду падающей волны, для некоторого фиксированного положения рупора записывается, очевидно, в виде
\begin{align} \label{eq:th:11}
	E = e^{-i jx} + \Gamma_k e^{i\varphi_k} e^{ihx} + \Gamma e^{i\varphi} e^{ihx}
\end{align}
где $x$ --- координата, отсчитываемая от конца волноводного тракта (см.~\cref{fig:4}); $h$ --- постоянная распространения волны в волноводе; $\Gamma_k e^{i \varphi_k}$ --- коэффициент отражения от конца тракта; $\Gamma e^{i \varphi}$ --- коэффициент отражения, обусловленный отражающим щитом. Выясните, на каком типе волны вы работаете и восстановите структуру электрического и магнитного полей в этой волне. Смещение антенны на величину $\Delta X$ приведет к появлению

в последнем члене в (11) дополнительного множителя $e^{ih_0 2\Delta X}$, связанного с дополнительным набегом фазы в свободном пространстве ($h_0 = \omega/c$ --- соответствующее волновое число). В результате для $|E|^2$ будем иметь

\begin{multline} \label{eq:th:12}
	|E|^2 = 1 + \Gamma_{\text{к}}^2 + \Gamma^2 + 2\Gamma_{k}\Gamma\cos(\varphi - \varphi_{k} + h_0 2\Delta X) +\\+ 2\Gamma_{k}\cos(2hx + \varphi_{k}) + 2\Gamma\cos(2hx + \varphi + h_0 2\Delta X).
\end{multline}

\begin{figure}[H]
	\centering
	\includegraphics[width=0.7\linewidth]{4}
	\caption{Волноводный тракт}
	\label{fig:4}
\end{figure}

Поскольку $\Gamma_{\text{к}}$ и $\Gamma$ достаточно малы, то квадратичными величинами в первом приближении можно пренебречь, т.е. опустить второй, третий и четвертый члены в (12); тогда эта формула упрощается к виду
\begin{align} \label{eq:th:13}
	|E|^2 \approx 1 + 2\Gamma_{k}\cos(2hx + \varphi_{k}) + 2\Gamma\cos(2hx + \varphi + k_0 2\Delta X).
\end{align}

Все возможные способы определения КНД на данной установке опираются в той или иной степени на эту формулу. Два таких способа изложены в задании к работе. Попытайтесь предложить другие способы. Еще раз подчеркнем, что КНД однозначно определяется коэффициентом $\Gamma$:
\begin{align} \label{eq:th:14}
	D = \frac{8\pi X}{\lambda} \Gamma.
\end{align}

